% !TEX root = ../main.tex
\section{Background}
\label{sec:background}

\subsection{Setting}

$N$ agents share a policy $\pi_{\bthe}$ with parameters $\bthe \in \Rb^d$ and
each interacts with its own copy of the same Markov Decision Process (MDP) $(\mathcal{S}, \mathcal{A},
P, r, \gamma)$. The objective is the usual expected discounted return
\begin{equation}
  J(\bthe) \;=\; \Exp_{\tau \sim \pi_{\bthe}}
  \left[\textstyle\sum_{t=0}^{\infty} \gamma^{t} r(s_t, a_t)\right],
\end{equation}
and the federation seeks $\bthe^\star = \arg\max_{\bthe} J(\bthe)$ using every
agent's samples while transmitting as little as possible.

\subsection{Policy gradient}

The policy gradient theorem \citep{williams1992reinforce, sutton2000pg}
expresses the gradient of $J$ as an expectation over the policy's own
trajectories,
\begin{equation}
  \nabla_{\bthe} J(\bthe) \;=\;
  \Exp_{\tau \sim \pi_{\bthe}}\!\left[\textstyle\sum_{t}
  \nabla_{\bthe} \log \pi_{\bthe}(a_t \mid s_t)\,
  A^{\pi_{\bthe}}(s_t, a_t)\right],
  \label{eq:pgt}
\end{equation}
where $A^{\pi}(s,a) = Q^{\pi}(s,a) - V^{\pi}(s)$ is the advantage of taking
$a$ in $s$. Each agent estimates \eqref{eq:pgt} from its own batch
$\mathcal{D}_i$ with the score-function estimator
\begin{equation}
  \bg_i \;=\; \frac{1}{|\mathcal{D}_i|}\sum_{\tau \in \mathcal{D}_i}\sum_{t}
  \nabla_{\bthe}\log\pi_{\bthe}(a_t \mid s_t)\,\hat{A}_t ,
  \label{eq:scorefn}
\end{equation}
which is line~7 of Algorithm~\ref{alg:fednpgadmm}. The advantages
$\hat{A}_t$ come from Generalized Advantage Estimation (GAE)
\citep{schulman2016gae}, an exponentially weighted sum of one-step temporal
difference errors of a learned value function $V$,
\begin{equation}
  \hat{A}_t \;=\; \sum_{l \ge 0} (\gamma\lambda)^l\, \delta_{t+l},
  \qquad
  \delta_t \;=\; r_t + \gamma V(s_{t+1}) - V(s_t),
\end{equation}
where $\lambda \in [0,1]$ trades bias for variance. We use $\lambda = 0.95$,
the value the paper states.

\subsection{Natural policy gradient}

Plain policy gradient ascends $\nabla_{\bthe} J$, which is sensitive to the
parameterization. An equal step in parameter space can mean a negligible or a
catastrophic change in the induced action distribution. NPG measures the step
in distribution space instead. To second order, the average Kullback--Leibler
(KL) divergence between the policy before and after a step $\Delta$ is a
quadratic form in the Fisher information matrix,
\begin{gather}
  \Exp_s\!\left[\KL\big(\pi_{\bthe}(\cdot\mid s)\,\big\|\,
                      \pi_{\bthe+\Delta}(\cdot\mid s)\big)\right]
  \;\approx\; \tfrac12\,\Delta^{\top}\bF(\bthe)\,\Delta,
  \label{eq:klquad}\\
  \bF(\bthe) \;=\;
  \Exp\!\left[\nabla_{\bthe} \log \pi_{\bthe}(a\mid s)\,
              \nabla_{\bthe} \log \pi_{\bthe}(a\mid s)^{\top}\right].
  \nonumber
\end{gather}
Maximizing the linearized objective $\bg^{\top}\Delta$ subject to
$\tfrac12\Delta^{\top}\bF\Delta \le \delta$ gives the natural-gradient step
\begin{equation}
  \bthe \;\leftarrow\; \bthe + \alpha\, \bF^{-1}\bg,
  \qquad
  \alpha \;=\; \sqrt{\frac{2\delta}{\bg^{\top} \bF^{-1} \bg}} ,
  \label{eq:npgstep}
\end{equation}
which is invariant to reparameterization and moves the policy by a KL
divergence of $\delta$ per step \citep{kakade2001npg, schulman2015trpo}.
Single-agent NPG, and Trust Region Policy Optimization (TRPO) with it, is
matrix-free: \eqref{eq:npgstep} needs only $\bF^{-1}\bg$, which conjugate
gradient computes from products $\bF v$, and each such product is a double
backward pass through the KL divergence. $\bF$ is never formed.

\subsection{The federated bottleneck}

Federating NPG naively means forming the global direction from every agent's
local gradient $\bg_i$ and local Fisher $\bH_i$:
\begin{equation}
  \bthe \leftarrow \bthe +
  \sqrt{\frac{2N\delta}
             {\big(\sum_{i} \bg_i\big)^{\!\top}
              \big(\sum_{i} \bH_i\big)^{-1} \sum_{i} \bg_i}}
  \left(\sum_{i=1}^{N} \bH_i\right)^{\!-1} \sum_{i=1}^{N} \bg_i .
  \label{eq:stdfednpg}
\end{equation}
This is what \citet{lan2023fednpgadmm} call \emph{standard-average FedNPG}.
The matrix-free trick no longer applies. The server needs each agent's $\bH_i$
to form the sum, and no agent can evaluate a Fisher-vector product on
another's data. Each agent must therefore uplink $\Oh(d^2)$ numbers per round.
Table~\ref{tab:commcost} shows what that means concretely for the three
environments studied here.

\begin{table}[t]
  \centering
  \caption{Per-agent uplink per iteration, at the paper's convention of 64 bits
  per transmitted parameter. $d$ is the policy parameter count for the
  architectures in Table~\ref{tab:table4}. The dense Fisher is reported for
  scale; it is never materialized in our implementation
  (Section~\ref{sec:baselines}).}
  \label{tab:commcost}
  \begin{tabular}{lrrrr}
    \toprule
    Environment & $d$ & Dense $\bH_i$ & $\Oh(d^2)$ / $\Oh(d)$ & Order \\
    \midrule
    Swimmer-v4  & \num{4868}   & \SI{0.095}{\giga\byte} & \num{2434}   & $10^{3.4}$ \\
    Hopper-v4   & \num{18438}  & \SI{1.36}{\giga\byte}  & \num{9219}   & $10^{4.0}$ \\
    Humanoid-v4 & \num{727074} & \SI{2114}{\giga\byte}  & \num{363537} & $10^{5.6}$ \\
    \bottomrule
  \end{tabular}
\end{table}
